\chapter{PENDAHULUAN}

\section{Latar Belakang}
Tugas 2 IF4020 Kriptografi meminta perancangan \textit{block cipher} baru,
implementasi enkripsi dan dekripsi dalam lima mode operasi, serta pengujian
\textit{avalanche effect}, entropi, dan histogram. Rancangan harus memiliki
substitusi berbasis tabel, permutasi, kunci putaran, dan operasi tambahan
di luar substitusi serta transposisi~\cite{spesifikasi}.

\Isi{Kembangkan motivasi rancangan, masalah yang ingin dijawab, dan alasan
memilih pendekatan Brainrot Cipher.}

\section{Tinjauan Cipher Pembanding}
AES menjadi pembanding yang relevan karena menggunakan blok 128 bit serta
transformasi substitusi dan permutasi~\cite{aes}. Implementasi Brainrot
menggunakan \textit{state} berukuran $4\times4$ byte, sepuluh putaran,
S-box hasil transformasi S-box AES, dan ekspansi kunci dengan pola serupa
AES-128. Parameter permutasi, matriks pencampuran, konstanta penambahan
kunci, dan substitusi dimodifikasi sebagaimana dijelaskan pada Bab III.

\Isi{Uraikan secara singkat persamaan, perbedaan, dan alasan pemilihan
komponen terhadap cipher pembanding. Jelaskan asal ide serta komponen
yang dirancang sendiri sesuai ketentuan tugas; jangan menyimpulkan
keamanan setara AES hanya dari kesamaan struktur.}

\section{Gagasan dan Ruang Lingkup}
\NamaAlgoritma{} menggunakan \textit{Substitution-Permutation Network}
(SPN), blok dan kunci utama sepanjang 128 bit, serta sepuluh putaran.
Program ditulis dalam Go dan menyediakan mode ECB, CBC, CFB, OFB, serta CTR.
Analisis memisahkan pengujian satu blok dari perilaku enkripsi pesan
utuh pada setiap mode.

\Isi{Lengkapi tujuan rancangan, alasan pemilihan parameter, dan batas
ruang lingkup laporan berdasarkan implementasi akhir.}
